\section{Introdução}
\label{sec:introducao}

A Atividade XII solicita que o grupo, com base no modelo de negócio desenvolvido até o momento, defina a estrutura do pitch, organizando os principais tópicos a serem apresentados, e crie a apresentação visual que o acompanhará, priorizando clareza, objetividade e apelo visual \cite{pucminas2026xii}. A orientação indica o uso de recursos visuais (imagens, gráficos, ícones), a economia de texto e o foco nas mensagens principais, e acolhe o uso de \textit{storytelling}. Exige, ainda, a identificação nominal dos integrantes da equipe.

Este documento cumpre a segunda entrega prevista, isto é, o texto que descreve a estrutura do pitch definida pelo grupo. A primeira entrega, a apresentação em formato \texttt{.pptx}, acompanha este texto (arquivo \texttt{Atividade\_XII\_Pitch\_PetCuida.pptx}). O conteúdo de ambos deriva das atividades anteriores: o protótipo e os testes com usuários (Atividade VI), os canais (Atividade VIII), os custos e a precificação (Atividade IX), a formalização (Atividade X) e a estrutura organizacional (Atividade XI).

Um pitch é, por definição, uma apresentação curta, destinada a despertar interesse e a conduzir o público a uma próxima conversa, e não a esgotar o plano de negócios. A literatura prática sobre apresentações para investidores recomenda poucas lâminas, tempo reduzido e foco em poucas mensagens \cite{kawasaki2004}. Por isso, a apresentação do grupo foi dimensionada para cerca de sete minutos e onze lâminas, e os detalhes numéricos permanecem nos documentos das atividades correspondentes.
